\section{Fixed standard activations with one nonlinear layer}
\label{sec:mixed-activation}

An exact extension is available for several standard nonlinear
activations if only the first hidden layer is nonlinear and the input
has no constant coordinate. These are explicit changes to the learner:
the result below does not concern a network with a nonlinear activation
at every hidden layer, and does not allow free first-layer biases.

Let $x\in\{-1,1\}^p$, let $N\ge2$, and write
\[
 Z_\theta(x)=vW^{L-1}\cdots W^2\phi(W^1x),\qquad L\ge2.
\]
The entries of $W^1$, $W^2,\ldots,W^{L-1}$, and $v$ have independent
centered Gaussian laws with variances $(N/p)\sigma^2$, $\sigma^2$,
and $\sigma^2/N$, respectively. Their quadratic energy is
\[
 \kappa_N(\theta)=\frac pN\|W^1\|_F^2+
 \sum_{j=2}^{L-1}\|W^j\|_F^2+N\|v\|_2^2.
\]
Condition on $|Z_\theta(x)|>1/2$ at every Boolean input,
and then on the observed labels.

Assume that $\phi$ is continuous and
\begin{equation}\label{eq:mixed-odd}
 \phi(t)-\phi(-t)=ct\quad(t\in\mathbb R),\qquad c>0,
\end{equation}
and assume in addition that either $t\phi(t)>0$ for $t\ne0$, or
$\phi$ is nondecreasing. GELU, SiLU, ReLU, and softplus satisfy these
conditions with $c=1$: here GELU is $t\Phi(t)$ with $\Phi$ the
standard normal distribution function, SiLU is $t/(1+e^{-t})$, ReLU is
$\max\{t,0\}$, and softplus is $\log(1+e^t)$. The first two preserve
the sign of nonzero inputs; the latter two are nondecreasing. A leaky ReLU with negative slope
$\alpha\in(0,1]$ satisfies them with $c=1+\alpha$.

\begin{lemma}[Exact energy for a paired history]
\label{lem:mixed-energy}
Let $H\subset\{-1,1\}^p\times\{-1,1\}$ be a nonempty history closed
under replacing $(x,y)$ by $(-x,-y)$. Assume the homogeneous linear margin constraints
are feasible. Let $A_*$ be the unique minimum-Euclidean-norm vector
satisfying $yA_*\cdot x\ge1/2$ on $H$. Suppose also that
$|A_*\cdot x|\ge1/2$ on the full cube. Then the minimum energy among
parameters with weak full-cube margin $1/2$ that fit $H$ is
\[
 E_*=L\left(\frac{4p\|A_*\|_2^2}{c^2}\right)^{1/L}.
\]
Every minimizing parameter represents $\operatorname{sign}(A_*\cdot x)$.
The actual conditioned Gaussian posterior therefore concentrates on
this sign function as $\sigma\downarrow0$, at each fixed $N,L$.
\end{lemma}
\begin{proof}
Put $b=vW^{L-1}\cdots W^2$. Equation~\eqref{eq:mixed-odd} gives
\[
 \frac{Z_\theta(x)-Z_\theta(-x)}2=A\cdot x,
 \qquad A=\frac c2 bW^1.
\]
The paired training constraints imply $yA\cdot x\ge1/2$, hence
$\|A\|_2\ge\|A_*\|_2$. Frobenius submultiplicativity and
arithmetic--geometric mean give
\begin{align*}
 \kappa_N(\theta)
 &\ge L\left(p\|W^1\|_F^2
       \prod_{j=2}^{L-1}\|W^j\|_F^2\|v\|_2^2\right)^{1/L}\\
 &\ge L\left(p\|W^1\|_F^2\|b\|_2^2\right)^{1/L}
 \ge L\left(4p\|A\|_2^2/c^2\right)^{1/L}\ge E_*.
\end{align*}

For attainment, let $u=A_*/\|A_*\|_2$ and put
\[
 e_1=(1,-1,0,\ldots,0)/\sqrt2.
\]
Choose unit vectors $e_2,\ldots,e_{L-1}\in\mathbb R^N$ and set
\[
 r=\left(4p\|A_*\|_2^2/c^2\right)^{1/(2L)},\qquad
 s_1=r\sqrt{N/p},\quad s_j=r\ (2\le j<L),\quad
 s_L=r/\sqrt N.
\]
Take $W^1=s_1e_1u^\top$, $W^j=s_je_je_{j-1}^\top$, and
$v=s_Le_{L-1}^\top$, with the evident interpretation when $L=2$.
The opposite first-layer pair cancels the even part of $\phi$, so
\[
 Z_\theta(x)=\frac c2\left(\prod_{j=1}^Ls_j\right)u\cdot x
             =A_*\cdot x.
\]
Each energy summand equals $r^2$, proving attainment.

At any energy minimizer all the preceding inequalities are equalities.
Thus $A=A_*$, and equality in
$\|bW^1\|_2\le\|b\|_2\|W^1\|_F$ forces
$W^1=\lambda b^\top u^\top$ for some $\lambda>0$. Consequently
\[
 Z_\theta(x)=F(u\cdot x),\qquad
 F(s)=\sum_i b_i\phi(\lambda b_is).
\]
If $t\phi(t)>0$ for $t\ne0$, each nonzero summand has the sign of
$s$, which proves the sign claim. If $\phi$ is nondecreasing, then
$F$ is nondecreasing, including when some $b_i$ are negative.
At least one training constraint for $A_*$ is active, since otherwise
$A_*$ could be scaled down. Its antipodal pair yields
\[
 F\left(\frac1{2\|A_*\|_2}\right)\ge\frac12,\qquad
 F\left(-\frac1{2\|A_*\|_2}\right)\le-\frac12.
\]
Every cube projection has absolute value at least
$1/(2\|A_*\|_2)$, so monotonicity again proves the sign claim.

There are finitely many Boolean sign functions, their weak-margin
parameter cells are closed, and the energy is coercive. All minimizers
have the same sign function, so every other consistent sign cell has a
strictly larger energy infimum. Increasing the readout in the attained
linear construction by a factor tending to one makes all margins
strict without changing the limiting energy. Lemma~\ref{lem:gaussian} now proves the posterior claim, including
any mixture of tied minimizers.
\end{proof}

\begin{theorem}[All widths for one nonlinear layer]\label{thm:mixed}
For every $k\ge3$, there is a Boolean target on $p=2k+2$ inputs,
computable by an $O(k)$-size bounded-fan-in Boolean circuit, and a fixed
sequence of distinct inputs forcing at least
$2^k-k-1\ge2^{k-1}$ mistakes in the strong-prior limit of the preceding
learner, at every fixed $N\ge2$ and $L\ge2$.
For each of the listed activations, the target's bounded-width-and-weight
complexity is exactly two at every depth $L\ge k+2$.
Thus the lower bound holds at every width at least that complexity.
For every $\eta>0$, all sufficiently small positive $\sigma$ make
posterior sampling wrong on every designated round with probability at
least $1-\eta$; posterior majority is wrong and the strict $2/3$
confidence criterion fails on all those rounds. The threshold may depend
on $N,k,L,\phi,\eta$; the input order does not depend on $N$ or $\sigma$.
\end{theorem}
\begin{proof}
Write $M=2^k-1$, add a variable Boolean coordinate $s$ to the selector
construction of Section~\ref{sec:selector}, and use the homogeneous target
\[
 F(s,x,y,z)=\operatorname{sign}\left(
   -Ms-(M+1)z+\sum_{i=0}^{k-1}2^i(x_i-y_i)\right).
\]
Equivalently, $F(s,x,y,z)=s f(sx,sy,sz)$ for the target in
\eqref{eq:target}. Coordinatewise multiplication by $s$ and the final
sign reversal cost only $O(k)$ Boolean gates.
The restriction $s=1$ is the comparator $f$, so the circuit size
is in fact $\Theta(k)$. Lift every history
point from Section~\ref{sec:selector} to $s=1$ and immediately follow it by its antipode, with the
opposite label. Lemma~\ref{lem:selector} gives, after the
paired history through query $t$, the unique unit-margin optimum
\[
 \widehat A_t=-ts-(t+1)z+\sum_i2^i(x_i-y_i).
\]
Its scores are odd nonzero integers on the full cube. Apply
Lemma~\ref{lem:mixed-energy} with $A_*=\widehat A_t/2$. The first point of each
next query pair, for $t=k,\ldots,M-1$, receives the incorrect positive
sign with posterior probability tending to one. The antipodal points
need not be counted. They are distinct and preserve a paired history
before each counted query. A finite-round union bound gives the
positive-scale sampled-prediction statement.

At depth $L=k+2$, take first-layer rows $\pm\widehat A_M/2^k$. In the first
linear hidden layer use the first row $(2/c,-2/c)$ and a zero second
row. Each of the remaining $k-1$ linear hidden layers multiplies the
first coordinate by two and keeps the second zero. Use readout
$(1,0)$. The resulting output is exactly $\widehat A_M$, with absolute margin at
least one. All coefficients are at most two for the listed activations,
and the width is two. Additional identity layers allow every
$L\ge k+2$.

Complexity one would give scalar layers with all coefficients bounded
by one. The odd projection, normalized from margin $1/2$ to unit
margin, would then have every coefficient bounded by $c\le2$.
Restricting its paired unit-margin constraints to $s=1$, the carry
inequalities in Lemma~\ref{lem:selector} instead force its highest-bit
coefficient difference divided by two to be at least
$2^{k-1}>2$. This contradiction proves that the complexity is exactly
two.
\end{proof}

Free first-layer biases remain a substantive obstacle. For example,
when $\beta\ge0$,
\[
 \operatorname{ReLU}(\beta+t)-\operatorname{ReLU}(\beta-t)
 =t+\operatorname{clip}(t,-\beta,\beta),
\]
so the odd projection is nonlinear. Additional nonlinear hidden layers
also destroy the linear odd-projection identity used above.
